\documentclass[11pt]{article}

\usepackage{amsmath,amssymb,amsthm,mathtools}
\usepackage[margin=1in]{geometry}
\usepackage{microtype}
\usepackage[hidelinks]{hyperref}

\newtheorem{theorem}{Theorem}[section]
\newtheorem{lemma}[theorem]{Lemma}
\newtheorem{corollary}[theorem]{Corollary}
\newtheorem{remark}[theorem]{Remark}

\newcommand{\dd}{\,\mathrm{d}}
\newcommand{\one}{\mathbf{1}}
\newcommand{\ip}[2]{\left\langle #1,#2\right\rangle}
\newcommand{\join}{\mathbin{\vee}}

\title{Two-Root Book-Edge Leaves and Their Cones:\\
A Geometric-Mean Graphon Bound}
\author{}
\date{}

\begin{document}

\maketitle

\begin{abstract}
For integers $m,a,b\geq1$, let $B_{m,a,b}$ be a book of $m$ triangles
sharing one edge, with $a$ pendant leaves at one end of that edge and $b$
pendant leaves at the other.  We prove, for every graphon $W$ of edge
density $p\geq1/2$,
\[
 t(B_{m,a,b},W)
 \geq p^{a+b+1}(2p-1)^m
 =\Phi_{B_{m,a,b}}(p).
\]
The proof replaces a previously unsuccessful signed-moment argument by a
geometric-mean compression under the edge-biased probability measure.  A
normalized-edge Cauchy--Schwarz estimate gives the exact mean needed for
Jensen's inequality.  We then prove the corresponding cone theorem for
$K_1\join B_{m,a,b}$ on the full interval $p\geq2/3$.  In the catalogue on
at most six vertices, the newly positive graphs are Atlas 35, 93, 112, and
180.
\end{abstract}


%%%%%%%%%%%%%%%%%%%%%%%%%%%%%%%%%%%%%%%%%%%%%%%%%%%%%%%%%%%%%%%%%%%%%%%%%%%%%%%
\section{Graphons and the book-edge families}
%%%%%%%%%%%%%%%%%%%%%%%%%%%%%%%%%%%%%%%%%%%%%%%%%%%%%%%%%%%%%%%%%%%%%%%%%%%%%%%

Let $(\Omega,\mu)$ be an arbitrary probability space.  A graphon is a
symmetric measurable function $W\colon\Omega^2\to[0,1]$.  For every finite
simple graph $F$, its ordinary, non-injective homomorphism density is
\[
 t(F,W)
 =
 \int_{\Omega^{V(F)}}
 \prod_{uv\in E(F)}W(x_u,x_v)
 \prod_{u\in V(F)}\dd\mu(x_u).
\]
All integrands below are measurable and bounded, so Tonelli--Fubini applies.
Write
\[
 p=t(K_2,W),
 \qquad
 d(x)=\int_\Omega W(x,y)\dd\mu(y).
\]

Fix an edge $rs$.  For integers $m,a,b\geq1$, define $B_{m,a,b}$ by
adjoining pairwise distinct page vertices $t_1,\ldots,t_m$, each adjacent
precisely to $r$ and $s$, then adjoining leaves $u_1,\ldots,u_a$ adjacent
only to $r$ and leaves $v_1,\ldots,v_b$ adjacent only to $s$.  Thus the
triangles $rst_i$ form an $m$-page book with spine $rs$, and both ends of
the spine receive a nonempty pendant-leaf set.  Interchanging $a$ and $b$
gives an isomorphic graph.  The earlier two-root triangle-leaf family is
the special case $m=1$.

The second family consists of the cones
\[
 C_{m,a,b}:=K_1\join B_{m,a,b}.
\]
Thus the new cone vertex is adjacent to every vertex of $B_{m,a,b}$, including
all of its pendant leaves.

Put $n=a+b$.  Once the spine has been properly coloured, every leaf has
$x-1$ independent colour choices, while every page vertex has $x-2$
choices after the spine has been coloured.  Hence
\begin{equation}
 \chi_{B_{m,a,b}}(x)=x(x-1)^{n+1}(x-2)^m,
 \qquad
 \chi(B_{m,a,b})=3.
 \label{eq:base-chromatic}
\end{equation}
For the cone, colouring the cone vertex first gives
\begin{equation}
 \chi_{C_{m,a,b}}(x)
 =x\chi_{B_{m,a,b}}(x-1)
 =x(x-1)(x-2)^{n+1}(x-3)^m,
 \qquad
 \chi(C_{m,a,b})=4.
 \label{eq:cone-chromatic}
\end{equation}

For $q=1-p>0$, direct substitution in the definition of the chromatic
target gives
\begin{align}
 \Phi_{B_{m,a,b}}(p)
 &=q^{n+m+2}\frac1q
   \left(\frac pq\right)^{n+1}
   \left(\frac{2p-1}{q}\right)^m
 =p^{n+1}(2p-1)^m,
 \label{eq:base-target}
 \\
 \Phi_{C_{m,a,b}}(p)
 &=q^{n+m+3}\frac1q\frac pq
   \left(\frac{2p-1}{q}\right)^{n+1}
   \left(\frac{3p-2}{q}\right)^m
 \notag\\
 &=p(2p-1)^{n+1}(3p-2)^m.
 \label{eq:cone-target}
\end{align}
Both right-hand sides are polynomials and therefore give the continued
values at $p=1$, namely $1$.


%%%%%%%%%%%%%%%%%%%%%%%%%%%%%%%%%%%%%%%%%%%%%%%%%%%%%%%%%%%%%%%%%%%%%%%%%%%%%%%
\section{The normalized-edge geometric-mean estimate}
%%%%%%%%%%%%%%%%%%%%%%%%%%%%%%%%%%%%%%%%%%%%%%%%%%%%%%%%%%%%%%%%%%%%%%%%%%%%%%%

The following lemma is the new analytic ingredient.

\begin{lemma}[Geometric mean of endpoint degrees]
\label{lem:edge-geometric-mean}
Suppose $p>0$, and put
\[
 Z(x,y):=\sqrt{d(x)d(y)}.
\]
Then
\begin{equation}
 \int_{\Omega^2}W(x,y)Z(x,y)\dd\mu(x)\dd\mu(y)\geq p^2.
 \label{eq:geometric-mean}
\end{equation}
Equivalently, under the edge-biased probability measure
\[
 \dd\rho(x,y)=\frac{W(x,y)}p\dd\mu(x)\dd\mu(y),
\]
one has $\mathbb E_\rho Z\geq p$.
\end{lemma}

\begin{proof}
If $d(x)=0$, then $W(x,y)=0$ for almost every $y$, because $W\geq0$ and
$d(x)=\int W(x,y)\dd\mu(y)$.  Thus all quotients below may be defined as
zero when one of their degree denominators vanishes.

Cauchy--Schwarz first gives
\begin{align*}
 \int_{\Omega^2}\frac{W(x,y)}{Z(x,y)}\dd\mu(x)\dd\mu(y)
 &=\int_{\Omega^2}
   \sqrt{\frac{W(x,y)}{d(x)}}
   \sqrt{\frac{W(x,y)}{d(y)}}
   \dd\mu(x)\dd\mu(y)
 \\
 &\leq
 \left(\int_{\Omega^2}\frac{W(x,y)}{d(x)}\dd\mu^2\right)^{1/2}
 \left(\int_{\Omega^2}\frac{W(x,y)}{d(y)}\dd\mu^2\right)^{1/2}
 \leq1.
\end{align*}
Indeed, the first integral in the product equals
$\mu\{x:d(x)>0\}\leq1$, and the second is identical after exchanging
$x$ and $y$.

Apply Cauchy--Schwarz a second time to
$\sqrt{WZ}$ and $\sqrt{W/Z}$.  The zero-degree observation justifies the
pointwise factorization on the support of $W$, and hence
\[
 p^2
 =\left(\int_{\Omega^2}W\dd\mu^2\right)^2
 \leq
 \left(\int_{\Omega^2}WZ\dd\mu^2\right)
 \left(\int_{\Omega^2}\frac WZ\dd\mu^2\right)
 \leq\int_{\Omega^2}WZ\dd\mu^2.
\]
This is \eqref{eq:geometric-mean}.
\end{proof}

\begin{lemma}[Convex book truncation]
\label{lem:convex-truncation}
For all integers $n\geq0$ and $m\geq1$, the function
\[
 f_{n,m}(z):=z^n\max\{2z-1,0\}^m,
 \qquad 0\leq z\leq1,
\]
where $z^0=1$, is convex and nondecreasing.
\end{lemma}

\begin{proof}
The function vanishes on $[0,1/2]$.  On $[1/2,1]$ it is a product of
$n$ copies of the nonnegative increasing affine factor $z$ and $m$ copies
of the nonnegative increasing affine factor $2z-1$.  Its first derivative
is a sum of products of nonnegative factors and positive slopes; its second
derivative is a sum over ordered pairs of differentiated factors, again
with only nonnegative factors and positive slopes.  It is therefore
nondecreasing and convex on $[1/2,1]$.  Its value at $1/2$ is zero and its
right derivative there is nonnegative (positive when $m=1$, zero when
$m\geq2$), whereas the left derivative is zero.  The one-sided derivatives
thus increase across the joining point, proving the assertion on all of
$[0,1]$.
\end{proof}


%%%%%%%%%%%%%%%%%%%%%%%%%%%%%%%%%%%%%%%%%%%%%%%%%%%%%%%%%%%%%%%%%%%%%%%%%%%%%%%
\section{Triangle books with leaves at both spine ends}
%%%%%%%%%%%%%%%%%%%%%%%%%%%%%%%%%%%%%%%%%%%%%%%%%%%%%%%%%%%%%%%%%%%%%%%%%%%%%%%

\begin{theorem}[Two-root book-edge leaf theorem]
\label{thm:two-root}
For every $m,a,b\geq1$ and every graphon $W$ of edge density
$p\in[1/2,1]$,
\[
 \boxed{
 t(B_{m,a,b},W)
 \geq p^{a+b+1}(2p-1)^m
 =\Phi_{B_{m,a,b}}(p).
 }
\]
This is the complete interval required by $\chi(B_{m,a,b})=3$.
\end{theorem}

\begin{proof}
The case $p=1$ is treated separately below, so assume first that
$1/2\leq p<1$.  Integrate out all leaves, and then independently integrate
out the $m$ page vertices.  With
\[
 c(x,y):=\int_\Omega W(x,z)W(y,z)\dd\mu(z),
\]
one obtains the exact identity
\begin{equation}
 t(B_{m,a,b},W)
 =\int_{\Omega^2}W(x,y)d(x)^a d(y)^b c(x,y)^m\dd\mu^2.
 \label{eq:two-root-identity}
\end{equation}
For numbers $A,B\in[0,1]$, one has $AB\geq A+B-1$.  Applying this to
$W(x,z)$ and $W(y,z)$ and integrating in $z$ gives, almost everywhere,
\[
 c(x,y)\geq d(x)+d(y)-1.
\]
Also $c\geq0$, and arithmetic--geometric mean gives
\[
 c(x,y)
 \geq\max\{d(x)+d(y)-1,0\}
 \geq\max\{2\sqrt{d(x)d(y)}-1,0\}.
 \label{eq:common-neighbour-lower}
\]

The factors $W$ and $c^m$ are symmetric in $x,y$.  Consequently,
\eqref{eq:two-root-identity} is unchanged if $a$ and $b$ are exchanged,
and it equals the average of the two resulting integrals.  Since all
factors are nonnegative, \eqref{eq:common-neighbour-lower} and the
arithmetic--geometric mean inequality imply, with $n=a+b$ and
$Z=\sqrt{d(x)d(y)}$,
\begin{align*}
 t(B_{m,a,b},W)
 &\geq\int_{\Omega^2}W(x,y)
 \frac{d(x)^a d(y)^b+d(x)^b d(y)^a}{2}
 \max\{2Z-1,0\}^m\dd\mu^2
 \\
 &\geq\int_{\Omega^2}W(x,y)
 (d(x)d(y))^{(a+b)/2}\max\{2Z-1,0\}^m\dd\mu^2
 \\
 &=p\,\mathbb E_\rho f_{n,m}(Z).
\end{align*}
Lemma~\ref{lem:convex-truncation}, Jensen's inequality,
Lemma~\ref{lem:edge-geometric-mean}, and monotonicity of $f_{n,m}$ now give
\[
 t(B_{m,a,b},W)
 \geq p f_{n,m}(\mathbb E_\rho Z)
 \geq p f_{n,m}(p)
 =p^{n+1}(2p-1)^m.
\]
The last expression is the target in \eqref{eq:base-target}.

If $p=1$, then $0\leq W\leq1$ and $\int W=1$ force $W=1$ almost
everywhere.  Hence $t(B_{m,a,b},W)=1$, which equals the polynomially
continued target.
\end{proof}

At the threshold $p=1/2$, the target is zero; no division by this factor or
by a triangle density occurs.  The proof above remains valid because the
edge-biased measure is well-defined.  Alternatively, the threshold statement
follows immediately from nonnegativity of homomorphism densities.


%%%%%%%%%%%%%%%%%%%%%%%%%%%%%%%%%%%%%%%%%%%%%%%%%%%%%%%%%%%%%%%%%%%%%%%%%%%%%%%
\section{A self-contained conditional cone lemma}
%%%%%%%%%%%%%%%%%%%%%%%%%%%%%%%%%%%%%%%%%%%%%%%%%%%%%%%%%%%%%%%%%%%%%%%%%%%%%%%

We reproduce the accepted link argument in the exact form needed for the
new family.  This makes clear why the next result is a conditional theorem,
not a general universal-vertex closure.

Define the rooted triangle density
\[
 \tau(x)=\int_{\Omega^2}
 W(x,y)W(x,z)W(y,z)\dd\mu(y)\dd\mu(z).
\]

\begin{lemma}[Weighted rooted-triangle inequality]
\label{lem:weighted-triangle}
If $p>0$, then for every integer $h\geq0$,
\[
 \int_\Omega d(x)^h\tau(x)\dd\mu(x)
 \geq\frac{2p-1}{p}\int_\Omega d(x)^{h+2}\dd\mu(x).
\]
\end{lemma}

\begin{proof}
Put $M_j=\int d^j$, let $T_W$ be the integral operator with kernel $W$,
and set
\[
 I_h=\int d^h\tau,
 \qquad
 J_h=\ip{d^h}{T_Wd}.
\]
The inequality $AB\geq A+B-1$, applied to $W(x,y),W(x,z)$ and
multiplied by $d(x)^hW(y,z)$, gives
\begin{equation}
 I_h\geq2J_h-pM_h.
 \label{eq:I-J}
\end{equation}
Put $U=1-W$, $u=T_U\one=1-d$, and $q=1-p$.  Since
\[
 T_Wd=d-q\one+T_Uu,
\]
we have
\begin{equation}
 J_h=M_{h+1}-qM_h+\ip{d^h}{T_Uu}.
 \label{eq:J-decomposition}
\end{equation}
By symmetry of $U$, twice the difference
\[
 \ip{d^h}{T_Uu}-\int d^hu^2
\]
equals
\[
 \int_{\Omega^2}U(x,y)
 \bigl(u(y)-u(x)\bigr)
 \bigl(d(x)^h-d(y)^h\bigr)\dd\mu^2\geq0,
\]
because $d=1-u$, so $u$ and $d^h$ are oppositely ordered.  Substitution
in \eqref{eq:J-decomposition} gives
\[
 J_h\geq M_{h+2}-M_{h+1}+pM_h.
\]
Together with \eqref{eq:I-J},
\[
 I_h\geq2M_{h+2}-2M_{h+1}+pM_h.
\]
Finally, subtracting $(2p-1)M_{h+2}/p$ from the right-hand side leaves
\[
 \frac1p\int_\Omega d(x)^h(d(x)-p)^2\dd\mu(x)\geq0.
\]
\end{proof}

For $d(x)>0$, define the probability measure
\[
 \dd\mu_x(y)=\frac{W(x,y)}{d(x)}\dd\mu(y),
\]
and regard $W$ as a graphon $W_x$ on $(\Omega,\mu_x)$.  Its edge density
is $z_x=\tau(x)/d(x)^2$.  Put $z_x=0$ when $d(x)=0$.  The preceding
lemma immediately yields, for every integer $N\geq2$,
\begin{equation}
 \int_\Omega d(x)^N z_x\dd\mu(x)
 \geq\left(2-\frac1p\right)\int_\Omega d(x)^N\dd\mu(x).
 \label{eq:weighted-link-average}
\end{equation}

\begin{lemma}[Conditional cone lemma]
\label{lem:conditional-cone}
Let $F$ have $N\geq2$ vertices, put $\phi(z)=\Phi_F(z)$, and suppose
that for some $\alpha\in[0,1)$:
\begin{enumerate}
\item every graphon $V$ of edge density $z\geq\alpha$ satisfies
$t(F,V)\geq\phi(z)$;
\item $\phi\geq0$, $\phi'\geq0$, and $\phi''\geq0$ on $[\alpha,1]$;
\item $\phi(\alpha)=0$.
\end{enumerate}
Then every graphon $W$ of edge density
\[
 p\geq\frac1{2-\alpha}
\]
satisfies
\[
 t(K_1\join F,W)\geq\Phi_{K_1\join F}(p).
\]
\end{lemma}

\begin{proof}
The endpoint $p=1$ is immediate, so assume $p<1$ and put
$c=2-1/p$.  The density hypothesis is exactly $c\geq\alpha$, and
$c\leq1$.  Convexity shows that every graphon $V$ of arbitrary edge
density $z\in[0,1]$ satisfies the global tangent bound
\[
 t(F,V)\geq\phi(c)+\phi'(c)(z-c).
\]
Indeed, for $z\geq\alpha$ this follows from the assumed graphon bound and
convexity.  For $z<\alpha$, the tangent is nonpositive: it is at most
$\phi(\alpha)=0$ at $\alpha$ and has nonnegative slope.  Nonnegativity of
homomorphism densities handles this second range.

Conditioning on the image of the cone vertex gives the exact identity
\[
 t(K_1\join F,W)=\int_\Omega d(x)^N t(F,W_x)\dd\mu(x).
\]
There is no division issue on $\{d=0\}$, because its contribution is zero.
Let $M_N=\int d^N$.  The tangent bound and
\eqref{eq:weighted-link-average} give
\begin{align*}
 t(K_1\join F,W)
 &\geq M_N\phi(c)
 +\phi'(c)\left(\int d(x)^Nz_x\dd\mu(x)-cM_N\right)
 \\
 &\geq M_N\phi(c)
 \geq p^N\phi(c),
\end{align*}
where the last step is Jensen and uses $\phi(c)\geq0$.

For $q=1-p>0$, one has $1-c=q/p$ and
$\chi_{K_1\join F}(x)=x\chi_F(x-1)$.  Therefore
\[
 p^N\phi(c)
 =q^N\chi_F\left(\frac pq\right)
 =q^{N+1}\chi_{K_1\join F}\left(\frac1q\right)
 =\Phi_{K_1\join F}(p).
\]
\end{proof}


%%%%%%%%%%%%%%%%%%%%%%%%%%%%%%%%%%%%%%%%%%%%%%%%%%%%%%%%%%%%%%%%%%%%%%%%%%%%%%%
\section{Cones over the two-root book family}
%%%%%%%%%%%%%%%%%%%%%%%%%%%%%%%%%%%%%%%%%%%%%%%%%%%%%%%%%%%%%%%%%%%%%%%%%%%%%%%

\begin{theorem}[Two-root book-edge leaf cones]
\label{thm:cones}
For every $m,a,b\geq1$ and every graphon $W$ of edge density
$p\in[2/3,1]$,
\[
 \boxed{
 t(C_{m,a,b},W)
 \geq p(2p-1)^{a+b+1}(3p-2)^m
 =\Phi_{C_{m,a,b}}(p).
 }
\]
This is the complete interval required by $\chi(C_{m,a,b})=4$.
\end{theorem}

\begin{proof}
By Theorem~\ref{thm:two-root}, the base $B_{m,a,b}$ has the graphon lower
bound
\[
 \phi_{n,m}(z)=z^{n+1}(2z-1)^m,
 \qquad n=a+b,
\]
throughout $z\geq1/2$.  On $[1/2,1]$, this polynomial is a product of
$n+1$ copies of the nonnegative increasing affine factor $z$ and $m$
copies of the nonnegative increasing affine factor $2z-1$.  Its first and
second derivatives are sums of products of nonnegative factors and
positive slopes.  Thus $\phi_{n,m}\geq0$, $\phi_{n,m}'\geq0$, and
$\phi_{n,m}''\geq0$ there, while $\phi_{n,m}(1/2)=0$.

Lemma~\ref{lem:conditional-cone}, with $\alpha=1/2$, applies for
\[
 p\geq\frac1{2-1/2}=\frac23.
\]
It proves the required graphon inequality.  Formula
\eqref{eq:cone-target} identifies the exact target.
\end{proof}


%%%%%%%%%%%%%%%%%%%%%%%%%%%%%%%%%%%%%%%%%%%%%%%%%%%%%%%%%%%%%%%%%%%%%%%%%%%%%%%
\section{Endpoints, equality examples, and catalogue consequences}
%%%%%%%%%%%%%%%%%%%%%%%%%%%%%%%%%%%%%%%%%%%%%%%%%%%%%%%%%%%%%%%%%%%%%%%%%%%%%%%

At $p=1/2$, the base target vanishes, and the balanced complete bipartite
graphon attains equality because $B_{m,a,b}$ contains a triangle.  At $p=2/3$,
the cone target vanishes, and the balanced complete tripartite graphon
attains equality because $C_{m,a,b}$ contains a $K_4$.  At $p=1$, the argument
given above shows that $W=1$ almost everywhere and both sides equal $1$.

More generally, for an integer $k\geq2$, the balanced complete $k$-partite
graphon $T_k$ has edge density $p=1-1/k$ and satisfies
\[
 t(F,T_k)=\frac{\chi_F(k)}{k^{v(F)}}=\Phi_F(p).
\]
This applies to $B_{m,a,b}$ for $k\geq2$, and to $C_{m,a,b}$ for $k\geq3$.
No claim is made here that these are all equality cases.

\begin{corollary}[New positive Atlas cases]
\label{cor:atlas}
The following formerly open catalogue rows satisfy the desired inequality
throughout their complete required intervals:
\[
\begin{array}{c|c|c|c|c|c|l}
\text{Atlas}&\text{graph6}&v&e&\chi&\text{structure}&\Phi_H(p)\\
\hline
35&\texttt{Dx\_}&5&5&3&B_{1,1,1}&p^3(2p-1)\\
93&\texttt{EC\char123 O}&6&6&3&B_{1,1,2}&p^4(2p-1)\\
112&\texttt{EhX\_}&6&7&3&B_{2,1,1}&p^3(2p-1)^2\\
180&\texttt{EjtW}&6&10&4&C_{1,1,1}&p(2p-1)^3(3p-2)
\end{array}
\]
Their chromatic polynomials are, respectively,
\[
 x(x-1)^3(x-2),
 \qquad
 x(x-1)^4(x-2),
 \qquad
 x(x-1)^3(x-2)^2,
 \qquad
 x(x-1)(x-2)^3(x-3).
\]
\end{corollary}

\begin{proof}
For unambiguous identification, the NetworkX Atlas edge lists are
\begin{align*}
35:\quad&
\{01,02,04,12,23\},\\
93:\quad&
\{03,04,14,24,34,35\},\\
112:\quad&
\{01,12,14,15,23,24,25\},\\
180:\quad&
\{01,04,12,13,14,15,23,34,35,45\}.
\end{align*}
In Atlas 35, $012$ is a triangle with one leaf at $0$ and one at $2$.
In Atlas 93, $034$ is a triangle with one leaf at $3$ and two at $4$.
In Atlas 112, the edge $12$ is shared by the triangles $124$ and $125$,
with one leaf at each end of $12$.  Thus Theorem~\ref{thm:two-root}
applies.  In Atlas 180, vertex $1$ is universal; deleting it leaves a copy
of $B_{1,1,1}$ on the other five
vertices, so Theorem~\ref{thm:cones} applies.  The orders, sizes, chromatic
polynomials, and targets follow from
\eqref{eq:base-chromatic}--\eqref{eq:cone-target}.
\end{proof}


%%%%%%%%%%%%%%%%%%%%%%%%%%%%%%%%%%%%%%%%%%%%%%%%%%%%%%%%%%%%%%%%%%%%%%%%%%%%%%%
\section{Repair of the earlier obstruction and hidden-assumption audit}
%%%%%%%%%%%%%%%%%%%%%%%%%%%%%%%%%%%%%%%%%%%%%%%%%%%%%%%%%%%%%%%%%%%%%%%%%%%%%%%

The earlier direct attempt for $B_{1,1,1}$ used
$W(x,z)W(y,z)\geq W(x,z)+W(y,z)-1$ and produced a difference of moments
with a negative coefficient.  Substituting a lower bound into that negative
term reverses the useful direction, so that proof cannot be repaired by
further algebra of the same two moments.

The proof here repairs precisely that step.  It keeps the positive part of
the common-neighbour estimate, symmetrizes the two rooted degree powers, and
compresses them to $Z=\sqrt{d(x)d(y)}$.  Lemma
\ref{lem:edge-geometric-mean} then supplies the missing exact lower bound
$\mathbb E_\rho Z\geq p$, and the convex function $f_{n,m}$ has no coefficient
with the wrong sign.

Theorems~\ref{thm:two-root} and \ref{thm:cones} apply directly to arbitrary
graphons on arbitrary probability spaces.  Quotients on zero-degree sets
are assigned zero only after proving that $W$ vanishes there almost
everywhere.  Every integral is of a bounded nonnegative measurable function,
apart from explicitly integrable normalized quotients.  The proof uses no
regularity, finite-step representation, finite-graph approximation,
injective-copy count, minimum-degree condition, pointwise lower bound on
$d$, or assumption that a conditioning event has positive mass.

The only standard analytic inputs are Tonelli--Fubini, arithmetic--geometric
mean, Cauchy--Schwarz, and Jensen.  The weighted link mechanism in
Section~4 is also the conditional cone lemma accepted in
\texttt{notes/six\_verified\_base\_cones.tex}; it was reproved here in full.


%%%%%%%%%%%%%%%%%%%%%%%%%%%%%%%%%%%%%%%%%%%%%%%%%%%%%%%%%%%%%%%%%%%%%%%%%%%%%%%
\section{Exact limitations}
%%%%%%%%%%%%%%%%%%%%%%%%%%%%%%%%%%%%%%%%%%%%%%%%%%%%%%%%%%%%%%%%%%%%%%%%%%%%%%%

The base theorem requires a book of triangles with one common spine edge,
and every remaining vertex must be a pendant leaf at one of the two spine
ends; both leaf sets must be nonempty in the catalogue predicate.  It does
not cover a leaf at a page vertex, a two-edge tail, a tree branching away
from the book, or attachments to unrelated core graphs.  The cone theorem
covers only $K_1\join B_{m,a,b}$ and depends on the explicitly verified
convex target $z^{a+b+1}(2z-1)^m$.  It is not a
general universal-vertex theorem and proves no arbitrary graph join.

No conclusion is asserted for graphs merely resembling these families, and
no numerical experiment or finite sampling is used in either proof.


%%%%%%%%%%%%%%%%%%%%%%%%%%%%%%%%%%%%%%%%%%%%%%%%%%%%%%%%%%%%%%%%%%%%%%%%%%%%%%%
\section{What the Lean formalization actually proves}
%%%%%%%%%%%%%%%%%%%%%%%%%%%%%%%%%%%%%%%%%%%%%%%%%%%%%%%%%%%%%%%%%%%%%%%%%%%%%%%

The machine-checked development is
\texttt{lean/Taeyoung/Methods/GeometricMean.lean} together with
\texttt{lean/Taeyoung/Methods/TwoRoot/\{Core,Rows\}.lean}, ending at
\texttt{satisfiesLowerBound\_\{35,93,112,180\}}.  It follows this note except at
one place.

\paragraph{The departure: Lemma~\ref{lem:convex-truncation} is not formalized.}
This note establishes that $f_{n,m}(z)=z^n\max\{2z-1,0\}^m$ is convex and
nondecreasing on $[0,1]$, by comparing one-sided derivatives across $z=1/2$, and
then applies Jensen.  The formalization never mentions convexity, derivatives,
or \texttt{ConvexOn}.  It instead proves the affine minorant that Jensen would
have produced, directly:
\begin{equation}
 f_{n,m}(p)+f_{n,m}'(p)(w-p)\leq f_{n,m}(w)
 \qquad\text{for all }w\geq0,\ p\geq1/2,
 \label{eq:lean-tangent}
\end{equation}
as \texttt{TwoRoot.tangent\_trunc}, where $f_{n,m}$ and $f_{n,m}'$ are the
formal product and product-rule derivative of the root list $1^n2^m$ in
\texttt{Methods/AffineProduct.lean}.  For $w\geq1/2$ the truncation is inactive
and \eqref{eq:lean-tangent} is the existing tangent lemma
\texttt{affineProd\_tangent} for products of nonnegative affine factors.  For
$w<1/2$ the right-hand side is zero, while the left-hand side is at most its
value at $w=1/2$, because the slope $f_{n,m}'(p)$ is nonnegative; and that value
is at most $f_{n,m}(1/2)=0$ by the same tangent lemma.  So the case split on the
sign of $2w-1$ happens once, pointwise in a lemma about real numbers, and never
inside an integral.

Integrating \eqref{eq:lean-tangent} against the edge-biased measure needs only
that the slope is nonnegative and that $\mathbb{E}_\rho[Z]\geq p$, which is
Lemma~\ref{lem:edge-geometric-mean}.  The result is the same bound this note obtains
from Jensen.  The two routes agree in strength; the formal one avoids
\texttt{ConvexOn}, one-sided derivatives, and the second-derivative argument in
the proof of Lemma~\ref{lem:convex-truncation}, none of which have a convenient
Mathlib interface for a truncated product.

\paragraph{Everything else transcribes directly.}
The geometric-mean estimate is
\texttt{GeometricMean.sq\_le\_integral\_geoIntegrand}, proved by the two
Cauchy--Schwarz steps of this note, carried out in $[0,\infty]$ so that the
unbounded quotients $W/d$ need no integrability side condition; the pointwise
factorizations there hold almost everywhere, and the null set is handled by the
same argument as in \texttt{notes/whiskering.tex}.  The two-vertex page bound
$H_0\geq d(x)+d(y)-1$ is \texttt{Link.le\_pageOp\_zero} and the
arithmetic--geometric mean step is \texttt{GeometricMean.two\_mul\_geoMean\_le}.
The symmetrization for $a\neq b$, needed only for Atlas 93, is
\texttt{MeasureTheory.integral\_prod\_swap} followed by
$d(x)d(y)^2+d(x)^2d(y)=Z^2(d(x)+d(y))\geq2Z^3$.  The cone (Atlas 180) uses the
already-formalized conditional cone lemma rather than the reproved version in
Section~4 of this note.

\end{document}
